\clearpage
% C057 | 三种主证明 + 两段角度/有向长度补充证明（保持原有证明完整）
\begin{proofbox}
  \textbf{\textcolor{CProof}{先证一个零门槛小工具：焦半径公式}}

  只证右焦点 $F(c,0)$（左焦点把图左右镜像即可）。双曲线上任取 $P(x,y)$，由
  $y^2=\dfrac{b^2}{a^2}x^2-b^2$、$c^2=a^2+b^2$ 得
  \[
    \begin{aligned}
      |PF|^2&=(x-c)^2+y^2\\
      &=\frac{c^2}{a^2}x^2-2cx+a^2
       =\left(\frac ca x-a\right)^2.
    \end{aligned}
  \]
  设 $e=c/a>1$、$p=b^2/a=ec-a$，于是
  \[
    \boxed{|PF|=|ex-a|}
    \quad\Longrightarrow\quad
    \begin{cases}
      |PF|=ex-a,&P\text{ 在右支},\\
      |PF|=a-ex,&P\text{ 在左支}.
    \end{cases}
  \]
  以下三法都只使用这一个工具（常规法再使用韦达）。

  \medskip
\end{proofbox}\clearpage\begin{proofbox}
  \textbf{\textcolor{CProof}{关键补充引理：焦点--准线成对；曲线上的点不分支}}

  \textbf{只需记一句：}焦点和同侧准线绑定，点在左支还是右支都可以代入。
  左、右两套焦点--准线为
  \[
    F_1(-c,0)\ \leftrightarrow\ \ell_1:x=-\frac{a^2}{c},\qquad
    F_2(c,0)\ \leftrightarrow\ \ell_2:x=\frac{a^2}{c}.
  \]
  对任意双曲线上的 $P(x,y)$，前页已经推得
  $|PF_2|=|ex-a|$。因为 $e=c/a$ 且 $a/e=a^2/c$，所以
  \[
    \boxed{|PF_2|=|ex-a|
     =e\left|x-\frac{a^2}{c}\right|
     =e\,d(P,\ell_2)}.
  \]
  以左焦点 $F_1(-c,0)$ 同法（或将 $x$ 换为 $-x$）：
  \[
    \boxed{|PF_1|=|ex+a|
     =e\left|x+\frac{a^2}{c}\right|
     =e\,d(P,\ell_1)}.
  \]
  因而对\textbf{每一个}双曲线上的点 $P$，\textbf{两套对应关系同时成立}：
  \[
    \boxed{\frac{|PF_1|}{d(P,\ell_1)}
      =\frac{|PF_2|}{d(P,\ell_2)}=e>1}.
  \]

  \textbf{为何左支仍能配右焦点？}
  若 $x\le -a$，则 $x<a^2/c$，右准线在 $P$ 右边，
  $d(P,\ell_2)=a^2/c-x$；于是
  \[
     |PF_2|=e\left(\frac{a^2}{c}-x\right)=a-ex>0.
  \]
  同理右支 $x\ge a$ 时，$|PF_2|=e(x-a^2/c)=ex-a$。
  \textbf{正负差别仅来自去绝对值，不是焦点--准线搭档发生了变化。}

  \textbf{与 C057 的联系：}即使焦点弦两点分居左右两支，
  既然直线通过 $F_2$，就可始终固定 $F_2$ 与 $\ell_2$；
  再根据共线关系消去角度即可。\textbf{切勿拿右焦点配左准线，亦勿把比值误写成 $1$。}
\end{proofbox}\clearpage\begin{proofbox}
  \textbf{\textcolor{CProof}{证明一（最快：方向余弦，一眼消去角度）}}

  \textbf{切入点：} 两个点都在\textbf{经过 $F$ 的同一条直线}上；以 $F$ 为起点，方向只影响 $\cos\theta$，目标是把它消掉。

  从 $F(c,0)$ 沿某射线走 $r=|FP|$，方向与 $x$ 正向夹角 $\theta$，则 $x=c+r\cos\theta$。记 $u=\cos\theta$。
  \begin{description}[style=nextline,leftmargin=2.8em,labelsep=0.5em,itemsep=0.4em]
    \item[\textbf{一步：写出射线上的倒数。}]
    若 $P$ 在右支，$r=e(c+ru)-a=p+eru$；若在左支，$r=a-e(c+ru)=-p-eru$。因此
    \[
      \boxed{\text{右支： }\frac1r=\frac{1-eu}{p}},\qquad
      \boxed{\text{左支： }\frac1r=\frac{-1-eu}{p}}.
    \]
    \item[\textbf{同支：相反方向，直接加。}]
    两点必在右支，且 $F$ 处于它们中间，两射线的余弦分别为 $u,-u$：
    \[
      \frac1{|AF|}+\frac1{|BF|}
      =\frac{1-eu}{p}+\frac{1+eu}{p}
      =\boxed{\frac2p=\frac{2a}{b^2}}.
    \]
    \item[\textbf{异支：同一方向，直接减。}]
    设 $A$ 为右支、$B$ 为左支，两点都在 $F$ 指向左侧的同一射线上，共用一个 $u$：
    \[
      \frac1{|AF|}-\frac1{|BF|}
      =\frac{1-eu}{p}-\frac{-1-eu}{p}
      =\boxed{\frac2p=\frac{2a}{b^2}}.
    \]
  \end{description}
  \textbf{优势：} 全程不设直线斜率、不解交点、不用韦达；还自动证明异支时右支上的点离右焦点更近。

  \medskip
\end{proofbox}\clearpage\begin{proofbox}
  \textbf{\textcolor{CProof}{证明二（最低门槛：只靠共线比例，无三角函数）}}

  \textbf{为什么设距离？} 题目要的是 $|AF|,|BF|$ 的倒数，所以\textbf{直接设} $r_A=|AF|,\ r_B=|BF|$；不设 $A(x_1,y_1),B(x_2,y_2)$ 四个未知量。焦半径公式立刻给出各点横坐标：
  \[
    \begin{array}{lll}
      \text{右支：}& r=ex-a,&\displaystyle x-c=\frac{r-p}{e};\\[4pt]
      \text{左支：}& r=a-ex,&\displaystyle x-c=\frac{-r-p}{e}.
    \end{array}\tag{*}
  \]
  这里仅用到 $ec-a=p$。又因为 $A,F,B$ 共线，\textbf{两条直角三角形的横边与斜边成比例}；同向取等号，反向取相反数。
  \begin{description}[style=nextline,leftmargin=2.8em,labelsep=0.5em,itemsep=0.45em]
    \item[\textbf{同支（$A-F-B$）：}]
    $A,B$ 都在右支，分别把 $(*)$ 代入“\textbf{方向相反}”的关系：
    \[
      \frac{x_A-c}{r_A}=-\frac{x_B-c}{r_B}
      \quad\Longrightarrow\quad
      \frac{r_A-p}{e r_A}=-\frac{r_B-p}{e r_B}.
    \]
    两边乘以 $e$ 并移项：
    \[
      1-\frac p{r_A}=-1+\frac p{r_B}
      \quad\Longrightarrow\quad
      \boxed{\frac1{r_A}+\frac1{r_B}=\frac2p}.
    \]
    \item[\textbf{异支（$F-A-B$）：}]
    $A$ 在右支、$B$ 在左支，使用“\textbf{方向相同}”的关系：
    \[
      \frac{x_A-c}{r_A}=\frac{x_B-c}{r_B}
      \quad\Longrightarrow\quad
      \frac{r_A-p}{e r_A}=\frac{-r_B-p}{e r_B}.
    \]
    同样只需移项：
    \[
      1-\frac p{r_A}=-1-\frac p{r_B}
      \quad\Longrightarrow\quad
      \boxed{\frac1{r_A}-\frac1{r_B}=\frac2p}.
    \]
  \end{description}
  \textbf{优势：} 只需要\textbf{设两段长度、焦半径公式、共线比例}三件事；不用记三角函数和二次方程。核心图像是\textbf{同支 $A-F-B$、异支 $F-A-B$}。

  \medskip
\end{proofbox}\clearpage\begin{proofbox}
  \textbf{\textcolor{CProof}{证明三（常规：直线联立 + 韦达，完整代数核验）}}

  设 $F(c,0)$。先取斜率存在的直线 $y=k(x-c)$，\textbf{且 $k\neq\pm b/a$}（否则仅有一个有限交点）。代入双曲线整理：
  \[
    (b^2-a^2k^2)x^2+2a^2ck^2x-a^2(c^2k^2+b^2)=0.
  \]
  两根 $x_1,x_2$ 对应两个有限交点，韦达给出
  \[
    x_1+x_2=\frac{2a^2ck^2}{a^2k^2-b^2},\qquad
    x_1x_2=\frac{a^2(c^2k^2+b^2)}{a^2k^2-b^2}.
  \]
  为避免一边写 $|ex_i-a|$ 一边反复拆绝对值，先计算\textbf{带符号的} $u_i=ex_i-a$：
  \[
    \begin{aligned}
      N:=u_1+u_2&=e(x_1+x_2)-2a
      =\frac{2ab^2(1+k^2)}{a^2k^2-b^2},\\
      M:=u_1u_2&=e^2x_1x_2-ae(x_1+x_2)+a^2
      =\frac{b^4(1+k^2)}{a^2k^2-b^2}.
    \end{aligned}
  \]
  因而（分母不为零）
  \[
    \boxed{\frac1{u_1}+\frac1{u_2}=\frac{u_1+u_2}{u_1u_2}
    =\frac NM=\frac{2a}{b^2}}.\tag{**}
  \]
  \begin{description}[style=nextline,leftmargin=2.8em,labelsep=0.5em,itemsep=0.4em]
    \item[\textbf{同支：}] 两点均在右支，$u_1,u_2>0$，于是 $r_A=u_1,r_B=u_2$。由 $(**)$ 得 $1/r_A+1/r_B=2a/b^2$。
    \item[\textbf{异支：}] 若 $A$ 在右支、$B$ 在左支，则 $u_A>0,u_B<0$；因此 $r_A=u_A,r_B=-u_B$，由 $(**)$ 得 $1/r_A-1/r_B=2a/b^2$，交换标签后写绝对值即可。
    \item[\textbf{斜率不存在：}] $x=c$ 时，$y=\pm b^2/a$，两交点均在右支；$r_A=r_B=b^2/a=p$，故 $1/r_A+1/r_B=2/p$。
    \item[\textbf{渐近线方向：}] $k=\pm b/a$ 时代入后的\textbf{二次项消失}，方程退化为一次方程，只得到一个有限交点；这两种情况不在定理的“两交点”前提内。
  \end{description}
  左焦点 $F(-c,0)$ 的结果由关于 $y$ 轴的对称性立即得到。

  \medskip
  \textbf{\textcolor{CProof}{三法选择}}

  \textbf{考场秒证：} 证明一；\quad
  \textbf{第一次理解、最低依赖：} 证明二；\quad
  \textbf{标准解析几何书面推导：} 证明三。
\end{proofbox}


\clearpage
\begin{proofbox}
  \textbf{\textcolor{CProof}{补充证明四（有向直线参数：一条公式统一同支和异支）}}

  \textbf{先定清两个不同的角。}本证明用 $\alpha$ 表示过右焦点 $F(c,0)$
  的\textbf{直线倾斜角}，取 $u=\cos\alpha$、$v=\sin\alpha$，
  $u^2+v^2=1$。对这条\textbf{整条直线}设
  \[
    P(t)=(c+tu,tv),\qquad t\in\mathbb R,\qquad |FP(t)|=|t|.
  \]
  这里 $t$ 是\textbf{有向位移}，不是距离；$t<0$ 表示沿方向 $(u,v)$ 的反向走。

  \textbf{直接联立，不用假设点在左支还是右支。}
  把 $P(t)$ 代入 $x^2/a^2-y^2/b^2=1$，乘以 $a^2b^2$，利用
  $c^2=a^2+b^2$、$u^2+v^2=1$，得到
  \[
  \begin{aligned}
    0&=b^2(c+tu)^2-a^2t^2v^2-a^2b^2\\
     &=(c^2u^2-a^2)t^2+2b^2cu\,t+b^4\\
     &=\bigl((cu-a)t+b^2\bigr)\bigl((cu+a)t+b^2\bigr).
  \end{aligned}
  \]
  因此，只要直线有\textbf{两个有限交点}（即 $cu\ne\pm a$），
  两个带符号的根为
  \[
    \boxed{t_1=\frac{b^2}{a-cu}},\qquad
    \boxed{t_2=-\frac{b^2}{a+cu}}.
  \]
  \textbf{关键：不去绝对值，直接对有向位移的倒数作差。}
  \[
    \boxed{\frac1{t_1}-\frac1{t_2}
      =\frac{a-cu}{b^2}+\frac{a+cu}{b^2}
      =\frac{2a}{b^2}}.
  \]
  下面才按位置\textbf{翻译成正距离}：
  \begin{description}[style=nextline,leftmargin=2.8em,labelsep=0.5em,itemsep=0.3em]
    \item[\textbf{同支（$|cu|<a$）：}]
    $t_1>0>t_2$，两交点在焦点两边，且都在右支。
    $r_1=t_1,r_2=-t_2$，因此
    $\displaystyle\frac1{r_1}+\frac1{r_2}=\frac1{t_1}-\frac1{t_2}=\frac{2a}{b^2}$。
    \item[\textbf{异支（$cu<-a$）：}]
    $0<t_1<t_2$，两交点位于同一正向射线上，
    先遇到右支、再遇到左支。
    $r_{\rm R}=t_1,r_{\rm L}=t_2$，故
    $\displaystyle\frac1{r_{\rm R}}-\frac1{r_{\rm L}}=\frac{2a}{b^2}$。
    \item[\textbf{异支（$cu>a$）：}]
    $t_1<t_2<0$，两交点都在反向射线上，先遇到右支、再遇到左支。
    $r_{\rm R}=-t_2,r_{\rm L}=-t_1$，故
    $\displaystyle\frac1{r_{\rm R}}-\frac1{r_{\rm L}}=\frac{2a}{b^2}$。
  \end{description}
  \textbf{边界说明：} $cu=\pm a$ 时有一个根的分母为零，
  即渐近线方向，只剩一个有限交点；原结论本来就要求两个有限交点。
  左焦点的版本由左右对称得到。
\end{proofbox}

\clearpage
\begin{proofbox}
  \textbf{\textcolor{CProof}{补充证明五（真正的射线焦半径式：为何分母必须为正）}}

  固定右焦点 $F(c,0)$，这次令 $\varphi$ 为射线 $\overrightarrow{FP}$
  的方向角，并且 $r=|FP|>0$。设 $w=\cos\varphi$，则
  $x=c+rw$、$p=b^2/a$、$e=c/a$。

  \begin{description}[style=nextline,leftmargin=2.8em,labelsep=0.5em,itemsep=0.35em]
    \item[\textbf{右支点：}]由 $r=ex-a$ 得
    \[
      r=e(c+rw)-a=p+erw
      \ \Longrightarrow\ 
      \boxed{r=\frac{p}{1-ew}=\frac{b^2}{a-cw}},
    \]
    这里\textbf{必须}有 $a-cw>0$ 才能使 $r>0$。
    \item[\textbf{左支点：}]由 $r=a-ex$ 得
    \[
      r=a-e(c+rw)=-p-erw
      \ \Longrightarrow\ 
      \boxed{r=-\frac{p}{1+ew}=-\frac{b^2}{a+cw}},
    \]
    这里\textbf{必须}有 $a+cw<0$，也就是 $w<-a/c$。
  \end{description}
  两个分母的符号限制来自\textbf{实际距离 $r>0$}，不是随便添加的附加条件。
  当焦点直线水平时，直线倾斜角 $\alpha=0$，
  但通向两个顶点的射线均朝左，故 $\varphi=\pi$、$w=-1$：
  \[
    r_{\rm R}=\frac{b^2}{a+c}=c-a,\qquad
    r_{\rm L}=-\frac{b^2}{a-c}=c+a.
  \]
  特别地 $a=3,b=4,c=5$，分别为 $2$、$8$。若向右选正向来定义
  \textbf{有向直线参数}，则得到 $t=-2,-8$；
  其中 $-8$ 可以描述\textbf{反向走 $8$}，却绝不能作为焦半径 $r$。
\end{proofbox}
